\documentclass[10pt]{beamer}
\usetheme{metropolis}
% all imports
\usepackage[utf8]{inputenc}
\usepackage[T1]{fontenc}
\usepackage{lmodern}
\usepackage{appendixnumberbeamer}
\usepackage{hyperref}
\usepackage{booktabs}
\usepackage{bm}
\usepackage[scale=2]{ccicons}
\usepackage[outputdir=build]{minted}
\usepackage{pgfplots}
\usepackage{array,colortbl,xcolor}
\usepgfplotslibrary{dateplot}
\usepackage{setspace}
\usepackage{etoolbox}
\usepackage{xspace}
\usepackage{tikz}
\usetikzlibrary{shapes,arrows,positioning,fit,backgrounds}
\usepackage{tkz-euclide}

\AtBeginEnvironment{quote}{\singlespacing}

\AtBeginEnvironment{quote}{\singlespacing}

% new commands
\newcommand{\themename}{\textbf{\textsc{metropolis}}\xspace}
\newcommand{\vect}[1]{\bm{#1}}
\newcommand{\myprime}[1]{{#1}^{\prime}}
\newcommand{\grad}[2]{\nabla_{#1} {#2}}
\newcommand{\dotp}[2]{{#1}^{\top}{#2}}
\newcommand{\dotpPright}[2]{{#1}^{\top}\left({#2}\right)}
\newcommand{\outerp}[2]{\left({#1}\right){#2}^{\top}}
\newcommand{\Jacobian}[2]{\frac{\partial #1}{\partial #2}}
\newcommand{\Vocab}{\mathbb{V}}
\DeclareMathOperator*{\argmin}{arg\,min}

% definitions
\definecolor{blue}{RGB}{159, 192, 176}
\definecolor{green}{RGB}{160, 227, 127}
\definecolor{orange}{RGB}{243, 188, 125}
\definecolor{red}{RGB}{253, 123, 84}
\definecolor{nephritis}{RGB}{39, 174, 96}
\definecolor{emerald}{RGB}{46, 204, 113}
\definecolor{turquoise}{RGB}{39, 174, 96}
\definecolor{green-sea}{RGB}{22, 160, 133}
% Tikzstyles for Computation Graphs

% nodes
\tikzstyle{noop} = [circle, draw=none, fill=red, minimum size = 10pt]
\tikzstyle{op} = [circle, draw=red, line width=1.5pt, fill=red!70, text=black, text centered, font=\bf \normalsize, minimum size = 25pt]
\tikzstyle{state} = [circle, draw=blue, line width=1.5pt, fill=blue!70, text=black, text centered, font=\bf \normalsize, minimum size = 25pt]
% \tikzstyle{gradient} = [circle, draw=green, line width=1.5pt, fill=green!60, text=black, text centered, font=\bf \normalsize, minimum size = 25pt]
\tikzstyle{gradient} = [circle, draw=nephritis, line width=1.5pt, fill=nephritis!60, text=black, text centered, font=\bf \normalsize, minimum size = 25pt]
\tikzstyle{textonly} = [draw=none, fill=none, text centered, font=\bf \normalsize]

% edges
% \tikzstyle{tedge}  = [draw, thick, >=stealth, ->]
\tikzstyle{tedge}  = [draw, thick, >=latex, ->]

% namedscope
\tikzstyle{namedscope} = [circle, draw=orange, line width=1.5pt, fill=orange!60, align=center, inner sep=0pt]

% \tikzstyle{container} = [draw=none, rectangle, dotted, inner ysep=1.5em]
% \tikzstyle{novertex} = [draw=none, fill=none, text centered]
% \tikzstyle{predicate} = [ellipse, draw, thick, text centered, rounded corners, minimum size=30pt]
% \tikzstyle{aux} = [rectangle, draw, thick, text centered, rounded corners, minimum size=30pt]
% \tikzstyle{ledge}  = [draw, dashed, thick, >=stealth, ->]
% \tikzstyle{pedge}  = [draw, thick, >=stealth, ->]



\title{IA na escola}
\date{\today}

\author{
  Felipe Salvatore\\
  \url{https://felipessalvatore.github.io/}
  \vspace{0.1 cm}
  \and\\ 
  Paula Moraes\\
  \url{https://www.linkedin.com/in/paula-moraes-137a2ab9/}\vspace{0.1 cm}
  \and\\ 
  Thiago Lira\\
  \url{https://github.com/ThiagoLira}\\\vspace{0.1 cm}
}

\institute{\textbf{IME-USP}: Institute of Mathematics and Statistics, University of São Paulo}

\titlegraphic{
            \hspace{9cm}
            \includegraphics[scale=0.2]{images/liamflogo.png}
}



\begin{document}
\nocite{DeepLearningbook}


\maketitle

\section{Introduction}
% definicao de IA em uma linha com exemplos 
% mudança de paradigma em IA. Track de IA do Bcc eh enviesado para logica http://bcc.ime.usp.br/principal/ia
% https://uspdigital.usp.br/jupiterweb/obterDisciplina?sgldis=MAC0425&verdis=3
% ML ser requisito ? Cada vez mais uso de dados e ML em IA, a disciplina de IA poderia ser a ponte entre IA classica e ML
% disciplina  de ML  https://uspdigital.usp.br/jupiterweb/obterDisciplina?sgldis=MAC0460&verdis=2
% aceitar cursos online de ML como requisito? cursos muito praticos FORMA MAO DE OBRA ESCRAVA
% NAO ESQUECER CRASE
% muda professores todo semestre com ementas levemente diferentes
%Ementa da Leliane vs. Ementa do Dennis
%logica pode ser optativa

% Leliane

% Busca Geral, Cega, Informada e Local
% Agentes Lógicos (I e II)
% Lógica de Primeira Ordem
% Cálculo de Situações
% Planejamento Automatizado I
% Planejamento Automatizado II - Planejador Progressivo, Regressivo e POP
% Planejamento Automatizado III - GRAPHPLAN e Heurísticas em Planejamento 
% Incerteza \& Raciocínio Probabilístico
% MDP (I e II)
% Aprendizado por Reforço (I e II)
% Aprendizado de Máquina: Redes Neurais 
% Aprendizado de Máquina: Árvore de Decisão
% Seminário de Deep Learning

% Dennis

% [introdução](aulas/intro.pdf) (cap 1 e 2)
% [busca cega](aulas/buscacega.pdf) (cap 3.1 a 3.4)
% [busca informada](aulas/buscainfo.pdf) (cap 3.5)
% [busca local](aulas/buscalocal.pdf) (cap 4)
% [busca competitiva](aulas/minimax.pdf) (cap 5)
% [agentes lógicos e lógica proposicional](aulas/logprop.pdf) (cap 7)
% [lógica de primeira ordem](aulas/lpo.pdf) (cap 8)
% [inferência lógica](aulas/inf.pdf) (cap 7 a 9)
% [incerteza](aulas/uncertainty.pdf) (cap 13)
% [raciocínio probabilístico I](aulas/independence.pdf) (cap 14)
% [raciocínio probabilístico II](aulas/bayesnets.pdf) (cap 14)
% [raciocínio probalístico temporal](aulas/dbn.pdf) (cap 15)
% [tomada de decisão simples](aulas/decision.pdf) (cap 16)
% [tomada de decisão sequencial](aulas/mdp.pdf) (cap 17)  
% [MDP: Algoritmos](aulas/mdp.pdf) (cap 17)  
% [POMDPs](aulas/pomdp.pdf) (cap 17)
% [aprendizagem por reforço](aulas/rl.pdf) (cap 21)
% [aprendizagem por reforço aproximada e regressão](aulas/rl2.pdf) (cap 18 e 21)
% [classificação](aulas/learning.pdf) (cap 18)
% 
%
%
%
%RL une os paradigmas classico e moderno de IA
%apresentar EP do pacman e um novo EP
%
% slide resumindo toda a nossa disrupcao e alteracoes na disciplina
%

\begin{frame}[fragile]{Basic idea}
\begin{itemize}
\item gen eh legal
\item limite eh dificil
\end{itemize}
\end{frame}

\begin{frame}[allowframebreaks]{References}

  \bibliography{my_references}
  \bibliographystyle{abbrv}

\end{frame}

\end{document}
